\documentclass[preprint,11pt,a4paper]{elsarticle}
\usepackage[T1]{fontenc}
\usepackage[utf8]{inputenc}
\usepackage[margin=25mm]{geometry}
\usepackage{amsmath,amssymb,booktabs,graphicx}
\usepackage[colorlinks=true,urlcolor=blue,citecolor=blue,linkcolor=black]{hyperref}
\usepackage{url}
\makeatletter
\def\ps@pprintTitle{\let\@oddhead\@empty\let\@evenhead\@empty
  \def\@oddfoot{\footnotesize Roland Cédric TAYO --- Teaching manuscript\hfill\thepage}
  \let\@evenfoot\@oddfoot}
\makeatother
\setlength{\parskip}{3pt}
\setlength{\emergencystretch}{2em}
\begin{document}

\begin{frontmatter}
\title{Replace with a precise title about your tested hypothesis}
\author{First Author, Second Author, Third Author, Fourth Author}
\address{INEM, Group XX, AI for Embedded Systems, 2026--2027}
\begin{abstract}
Replace this text with 120--180 words: problem, hypothesis, method, exact board, measured result and main limitation. Do not invent a result before the experiment. The final PDF must contain no more than 15 pages including references and appendices.
\end{abstract}
\begin{keyword}
Embedded inference \sep forecasting \sep replace with your optimisation
\end{keyword}
\end{frontmatter}
\section{Introduction}
Explain the local prediction need and the board constraint. State one testable hypothesis.
\section{Related work}
Describe the original TSMixer idea~\cite{chen2023tsmixer} and its difference from the reduced TD network. Cite an optimisation source and the target documentation. The target is a Raspberry Pi running Linux. Any optimised-library comparator must be separate from the handwritten implementation.
\section{Method}
Explain data generation, split, training-only normalisation, network and baselines. Define each symbol before using it. Show the implemented change with a diagram or short pseudocode.
\section{Experimental protocol}
Give the Pi model, OS architecture, CPU frequency, temperature, load, compiler, flags, versions, checksums, correctness tolerance, repetitions and timing scope. Explain dispersion and distinguish explicit buffer bytes from process memory.
\section{Results}
Insert actual training/validation curves, observed-versus-predicted traces and board reference-versus-variant measurements. Include raw-data paths and the commit for every campaign. Do not reuse the separate example trace as your TD dataset or result.
\section{Discussion and limitations}
Explain where the change helps or fails. Distinguish observations from hypotheses. Discuss the synthetic-data limitation and the whole-network effect.
\section{Conclusion and next experiment}
Answer the hypothesis and propose one next test with a measurable decision criterion.
\section*{Reproducibility and contributions}
Add repository URL, submission tag, commands, manifests, member contributions, reused code and substantial assistant usage. The teacher account to invite on GitHub is \texttt{rolln7drktayau}.
\bibliographystyle{elsarticle-num}
\bibliography{references}
\end{document}
